%! Exponential Growth and Decay — Unit 6, Lesson 6.1 — Exit Ticket (Answer Key)
\documentclass[10pt]{article}
\usepackage{algebra2-article}
\usepackage{algebra2-key}   % pulls in -boxes; do NOT also load -boxes
\newcommand{\TallMath}[1]{$\displaystyle #1\rule[-1.4em]{0pt}{3.2em}$}

\begin{document}

\pageheader{Unit 6, Lesson 6.1}{Exit Ticket --- Answer Key}
\namedateperiod

\vspace{0.08in}

No notes. Show the two numbers ($a$ and $b$) for every model you build.

\begin{enumerate}[leftmargin=1.8em, itemsep=9pt]
  \item \textbf{Percent to factor.} Write the base $b$ for each description.

        \vspace{3pt}
        A quantity \textbf{grows} $7\%$ each year: \; $b=$ \ans{1.07} \qquad
        A quantity \textbf{loses} $30\%$ each year: \; $b=$ \ans{0.70}

  \item \textbf{Build it, then use it.} A rare-book collection is worth \$$2400$ today and gains
        $7\%$ in value each year.

        \vspace{3pt}
        $a=$ \ans{2400} \quad $b=$ \ans{1.07} \quad
        $V(t)=$ \ans{$2400(1.07)^{t}$}

        \vspace{4pt}
        $V(3)=$ \ans{$2400(1.225043)$} $\approx$ \ans{\$$2940.10$}

        \vspace{3pt}
        Write one sentence saying what $V(3)$ means about the collection.
        \ansline{If the collection keeps gaining $7\%$ a year, it will be worth about \$$2940.10$
        three years from now --- roughly \$$540$ more than today.}

  \item \textbf{From a table.} Write the model in the form $y=ab^{x}$.
    \begin{center}
    \renewcommand{\arraystretch}{1.3}
    \begin{tabular}{|c|c|c|c|c|}
    \hline
    $x$ & $0$ & $1$ & $2$ & $3$ \\ \hline
    $y$ & $5$ & $20$ & $80$ & $320$ \\ \hline
    \end{tabular}
    \end{center}
    $a=$ \ans{5} \quad $b=$ \ans{4} \quad $y=$ \ans{$5\cdot 4^{x}$}

  \item \textbf{SOL-style multiple choice.} A town's population is $4000$ and is decreasing by $6\%$
        each year. Which model gives the population $P$ after $t$ years?
        \begin{enumerate}[label=\Alph*., leftmargin=1.9em, itemsep=1pt]
          \item $P(t) = 4000(0.94)^{t}$ \quad \textcolor{keyred}{\textbf{$\leftarrow$ correct}}
          \item $P(t) = 4000(1.06)^{t}$
          \item $P(t) = 4000(0.06)^{t}$
          \item $P(t) = 4000 - 0.06t$
        \end{enumerate}
\end{enumerate}

\begin{teachernote}
\textbf{Answer 4: A.} Each distractor is one specific error. \textbf{B} adds the rate instead of
subtracting it ($1+r$ for a decrease) --- the population \emph{grows} $6\%$ a year. \textbf{C} is the
percent-as-base error: only $6\%$ of the town survives each year, a $94\%$ collapse. \textbf{D} is
the family error --- a constant \emph{amount} subtracted each year is linear, and it loses
$0.06$ of a person annually, which should be visibly absurd.

\textbf{Sort the collected tickets into four piles} to decide how Lesson 6.2 opens:
\emph{got it} / \emph{sign error} (item 1's second answer written $1.30$, or item 4 answered B) /
\emph{percent as base} (item 1 written $0.07$, or item 4 answered C) / \emph{linear reflex}
(item 4 answered D). The middle two piles are both ``the base is not the percent'' and are worth a
three-minute re-teach off the Warm-Up jacket: after $+7\%$ you have $107\%$; after $-30\%$ you have
$70\%$.

\textbf{Item 3 is the low bar} --- a table hands over both numbers with no translation at all.
Anyone who misses it needs the constant-ratio idea itself retaught, not the percent rule, and will
struggle in 6.5 where the same reading drives the regression.

\textbf{Item 2's sentence} is the \textbf{A2.F.2f} evidence of the day (evaluate and interpret in
context). Accept any answer that names the dollar amount \emph{and} ties it to three years out; do
not accept ``$2940.10$'' alone.
\end{teachernote}

\end{document}
